% Part: first-order-logic
% Chapter: beyond
% Section: second-order-logic

\documentclass[../../../include/open-logic-section]{subfiles}

\begin{document}

\olfileid{fol}{byd}{sol}

\olsection{द्वितीय-क्रम तर्कशास्त्र}

द्वितीय-क्रम तर्कशास्त्राची भाषा व्यक्तींच्या !!a{domain} वरच नव्हे, तर त्या
!!{domain} वरील संबंधांवरही संख्यापन करण्याची मुभा देते. प्रथम-क्रम
भाषा~$\Lang{L}$ दिली असता, प्रत्येक $k$ साठी $k$-स्थानी संबंधांवर मान घेणारी
$R$ ही !!{variable}s जोडली जातात आणि त्या चरांवर संख्यापन करण्याची मुभा दिली
जाते. $R$ हे एखाद्या $k$-स्थानी संबंधाचे !!a{variable} असेल आणि $t_1$,
\dots,~$t_k$ ही नेहमीची (प्रथम-क्रम) पदे असतील, तर
$\Atom{R}{t_1,\dots,t_k}$ हे अण्वीय !!{formula} आहे. याखेरीज !!{formula}s चा
संच प्रथम-क्रम तर्कशास्त्राप्रमाणेच, फक्त द्वितीय-क्रम संख्यापनासाठी अतिरिक्त
उपनियम घालून, परिभाषित केला जातो. लक्षात घ्या की आपल्याकडे !!{identity} फक्त
प्रथम-क्रम पदांसाठी आहे: $R$ आणि $S$ ही समान स्थानसंख्या~$k$ असलेल्या
संबंधांची !!{variable}s असतील, तर $\eq[R][S]$ हे पुढील सूत्राचे संक्षिप्त रूप
म्हणून परिभाषित करता येते:
\[
\lforall[x_1 \dots][\lforall[x_k][(\Atom{R}{x_1, \dots, x_k} \liff
  \Atom{S}{x_1, \dots, x_k})]].
\]

द्वितीय-क्रम तर्कशास्त्राचे नियम संख्यापकांचे नियम नव्या द्वितीय-क्रम चरांपर्यंत
विस्तारित करतात. मात्र ही चरे $\Lang{L}$ मधील !!{predicate}s शी आणि अधिक
सामान्यपणे $\Lang{L}$ मधील !!{formula}s शी कशी परस्परक्रिया करतात, हे येथे
थोडे काळजीपूर्वक स्पष्ट करावे लागते. किमान चौकटीत संबंधचरे पदे म्हणून गणली
जातात; म्हणून पुढील प्रकारच्या अनुमानांना मुभा असते:
\[
!A(R) \Proves \lexists[R][!A(R)]
\]
पण $\Lang{L}$ ही अंकगणिताची भाषा असून $<$ हे स्थिर संबंधचिन्ह असेल, तर पुढील
अनुमानही ग्राह्य असावे अशी अपेक्षा असते:
\[
x < y \Proves \lexists[R][\Atom{R}{x,y}]
\]
किंवा दिलेल्या !!{formula}~$!A$ साठी,
\[
!A(x_1, \dots, x_k) \Proves \lexists[R][\Atom{R}{x_1,\dots,x_k}]
\]
आणखी सामान्यपणे, पुढील प्रकारच्या अनुमानांना मुभा द्यावीशी वाटू शकते:
\[
\Subst{!A}{\lambd[\vec x][!B(\vec x)]}{R} \Proves
\lexists[R][!A]
\]
येथे $\Subst{!A}{\lambd[\vec x][!B(\vec x)]}{R}$ म्हणजे $!A$ मधील
$\Obj{R}{t_1,\dots,t_k}$ या रूपाचे प्रत्येक अण्वीय !!{formula}
$!B(t_1, \dots, t_k)$ ने बदलल्यावर मिळणारे सूत्र. हा शेवटचा नियम
{\em गुणधर्माधारित संबंधरचना-रूपबंध} असण्याशी समतुल्य आहे; म्हणजे
\[
\lexists[R][\lforall[x_1, \dots, x_k][(!A(x_1, \dots, x_k) \liff
\Atom{R}{x_1, \dots, x_k})]],
\]
या रूपाचे, द्वितीय-क्रम भाषेतील प्रत्येक !!{formula}~$!A$ साठी एक,
स्वयंसिद्धक असते; त्यात $R$ हे मुक्त चर नसते. (सराव: $R$ ला $!A$ मध्ये येऊ
दिल्यास हा रूपबंध विसंगत होतो, हे दाखवा!)

तर्कशास्त्रज्ञ “द्वितीय-क्रम तर्कशास्त्राची स्वयंसिद्धके” म्हणतात तेव्हा
सहसा प्रथम-क्रम तर्कशास्त्राचा द्वितीय-क्रम संख्यापक-नियमांनी केलेला किमान
विस्तार आणि गुणधर्माधारित संबंधरचना-रूपबंध त्यांना अभिप्रेत असतो. पण या
स्वयंसिद्धकांच्या आणि नियमांच्या दुर्बल उपप्रणालींचा अभ्यास करणेही अनेकदा
उपयुक्त असते. उदाहरणार्थ, गुणधर्माधारित संबंधरचनेचा स्वयंसिद्धक-रूपबंध पूर्ण
सामान्यतेत \emph{स्वसमावेशक} आहे: द्वितीय-क्रम संख्यापक असलेल्या
!!a{formula} ने “परिभाषित” होणारा $\Atom{R}{x_1, \dots, x_k}$ हा संबंध
अस्तित्वात आहे असे विधान करण्याची तो मुभा देतो; आणि हे संख्यापक अशा सर्व
संबंधांच्या संचावर मान घेतात---ज्या संचात $R$ स्वतःही येतो!{} विसाव्या
शतकाच्या सुमारास रसेलच्या विरोधापत्तीला दिलेली एक सामान्य प्रतिक्रिया म्हणजे
अशा परिभाषांनाच तिच्यासाठी दोषी धरणे आणि गणिताचा पाया विकसित करताना त्या
टाळणे. !!{formula}~$!A$ मध्ये द्वितीय-क्रम संख्यापक वापरण्यास मनाई केली, तर
गुणधर्माधारित संबंधरचनेचे काहीसे दुर्बल असे {\em स्वसमावेशरहित} रूप मिळते.

चिन्हार्थाच्या दृष्टीने द्वितीय-क्रम !!{structure} मध्ये त्या भाषेसाठीची एक
प्रथम-क्रम !!{structure} आणि ज्या संबंधांवर द्वितीय-क्रम संख्यापक मान घेतात
असा त्या !!{domain} वरील संबंधांचा संच असतो (अधिक नेमके सांगायचे, तर प्रत्येक
$k$ साठी स्थानसंख्या~$k$ असलेल्या संबंधांचा एक संच असतो). अर्थात,
!!{derivation} प्रणालीत गुणधर्माधारित संबंधरचना समाविष्ट असेल, तर
“द्वितीय-क्रम भागात” ती स्वयंसिद्धके पूर्ण करण्याइतके संबंध असावेत ही
अतिरिक्त अट येते---नाहीतर !!{derivation} प्रणाली निर्दोष राहत नाही!{} पुरेसे
संबंध उपलब्ध असल्याची खात्री करण्याचा एक सोपा मार्ग म्हणजे द्वितीय-क्रम भागात
प्रथम-क्रम भागावरील \emph{सर्व} संबंध घेणे. अशा !!a{structure} ला
\emph{पूर्ण} म्हणतात आणि एका अर्थाने तीच भाषेची “अभिप्रेत !!{structure}”
असते. आपण फक्त पूर्ण !!{structure}s चा विचार केला, तर त्याला
\emph{पूर्ण} द्वितीय-क्रम चिन्हार्थमीमांसा म्हणतात. अशा वेळी !!{structure}
निर्दिष्ट करणे म्हणजे फक्त तिचा प्रथम-क्रम भाग निर्दिष्ट करणे; कारण
द्वितीय-क्रम भागातील घटक त्यावरून अप्रत्यक्षपणे ठरतात.

सारांशतः, द्वितीय-क्रम तर्कशास्त्राविषयी बोलताना काही संदिग्धता असते.
!!{derivation} प्रणालीच्या दृष्टीने पुढीलपैकी कोणतीही गोष्ट अभिप्रेत असू शकते:
\begin{enumerate}
\item काही गुणधर्माधारित संबंधरचना-स्वयंसिद्धकांसह “किमान” द्वितीय-क्रम
  !!{derivation} प्रणाली.
\item पूर्ण गुणधर्माधारित संबंधरचना असलेली “मानक” द्वितीय-क्रम
  !!{derivation} प्रणाली.
\end{enumerate}
चिन्हार्थाच्या दृष्टीने पुढीलपैकी कशातही स्वारस्य असू शकते:
\begin{enumerate}
\item “दुर्बल” चिन्हार्थमीमांसा, जिच्यात !!a{structure} मध्ये प्रथम-क्रम
  भाग आणि गुणधर्माधारित संबंधरचना-स्वयंसिद्धके पूर्ण करण्याइतका मोठा
  द्वितीय-क्रम भाग असतो.
\item “मानक” द्वितीय-क्रम चिन्हार्थमीमांसा, जिच्यात फक्त पूर्ण
  !!{structure}s चा विचार केला जातो.
\end{enumerate}
तर्कशास्त्रज्ञांनी कोणती !!{derivation} प्रणाली किंवा कोणती चिन्हार्थमीमांसा
अभिप्रेत आहे हे सांगितले नाही, तर सहसा प्रत्येक यादीतील दुसरी बाब त्यांना
अभिप्रेत असते. पुढे दिसेल त्याप्रमाणे, ही चिन्हार्थमीमांसा वापरण्याचा फायदा असा
की तिच्यात अनेक स्वाभाविक गणितीय रचनांची समरूपतेपर्यंत एकमेव वर्णने मिळतात;
त्याच वेळी !!{derivation} प्रणाली बरीच समर्थ असून या चिन्हार्थमीमांसेसाठी
निर्दोष असते. मात्र !!{derivation} प्रणाली या चिन्हार्थमीमांसेसाठी
\emph{संपूर्ण नाही}; इतकेच नव्हे, तर प्रभावी रीतीने दिलेली \emph{कोणतीही}
!!{derivation} प्रणाली पूर्ण द्वितीय-क्रम चिन्हार्थमीमांसेसाठी संपूर्ण नसते.
दुसरीकडे, दुर्बल चिन्हार्थमीमांसेसाठी !!{derivation} प्रणाली \emph{संपूर्ण
आहे}, हे आपण पाहू; म्हणून एखादे वाक्य सिद्ध करता येत नसेल, तर ते असत्य असणारी
\emph{एखादी तरी} !!{structure} असते, जरी ती पूर्ण रचना असेलच असे नाही.

द्वितीय-क्रम तर्कशास्त्राची भाषा बरीच अभिव्यक्तिसमर्थ आहे. एकस्थानी संबंधांची
!!{domain} च्या उपसंचांशी ओळख करता येते; म्हणून विशेषतः या संचांवर संख्यापनही
करता येते. उदाहरणार्थ, स्वाभाविक संख्यांसाठीचे आगमन एकाच स्वयंसिद्धकाने
व्यक्त करता येते:
\[
\lforall[R][((\Atom{R}{\Obj{0}} \land \lforall[x][(\Atom{R}{x} \lif
    \Atom{R}{x'})]) \lif \lforall[x][\Atom{R}{x}])].
\]
अंकगणिताच्या भाषेत $\Obj 0, \Obj \prime, +, \times$ आणि $<$ ही चिन्हे घेतली,
तर त्यांचे वर्तन वर्णिण्यासाठी पुढील स्वयंसिद्धके जोडता येतात:
\begin{enumerate}
\item $\lforall[x][\lnot x' = \Obj 0]$
\item $\lforall[x][\lforall[y][(s(x) = s(y) \lif x = y)]]$
\item $\lforall[x][(x + \Obj 0) = x]$
\item $\lforall[x][\lforall[y][(x + y') = (x + y)']]$
\item $\lforall[x][(x \times \Obj 0) = \Obj 0]$
\item $\lforall[x][\lforall[y][(x \times y') = ((x \times y) + x)]]$
\item $\lforall[x][\lforall[y][(x < y \liff \lexists[z][y = (x + z')])]]$
\end{enumerate}
आपण पूर्ण द्वितीय-क्रम चिन्हार्थमीमांसा वापरत असू, तर ही स्वयंसिद्धके आणि
वरील आगमनाचे स्वयंसिद्धक मिळून अंकगणिताचे मानक प्रतिमान असलेल्या
!!{structure}~$\Struct{N}$ चे समरूपतेपर्यंत एकमेव वर्णन देतात, हे दाखवणे
अवघड नाही. ही स्वयंसिद्धके सत्य असलेली कोणतीही !!{structure}~$\Struct{M}$
दिली असता, $\Nat$ वरील नेहमीचे पुनरावर्तन वापरून $\Nat$ पासून
$\Struct{M}$ च्या !!{domain} कडे जाणारे $f$ हे फलन असे परिभाषित करा की
$f(0) = \Assign{\Obj 0}{M}$ आणि $f(x+1) = \Assign{\prime}{M}(f(x))$.
$\Nat$ वरील नेहमीचे आगमन आणि $\Struct M$ मध्ये स्वयंसिद्धके (1) आणि~(2)
सत्य असण्याचा उपयोग केल्यावर $f$ हे !!{injective} आहे असे दिसते. $f$ हे
!!{surjective} आहे हे पाहण्यासाठी, $f$ च्या व्याप्तीत असलेल्या
$\Domain{M}$ मधील घटकांचा संच $P$ घ्या. $\Struct M$ पूर्ण असल्यामुळे $P$
हा द्वितीय-क्रम !!{domain} मध्ये असतो. $f$ च्या रचनेवरून
$\Assign{\Obj 0}{M}$ हा $P$ मध्ये असतो आणि $P$ हा
$\Assign{\prime}{M}$ खाली बंद असतो. $\Struct M$ मध्ये आगमनाचे स्वयंसिद्धक
सत्य असणे (विशेषतः $P$ साठी) $P$ हा $\Struct M$ च्या संपूर्ण प्रथम-क्रम
!!{domain} इतकाच असल्याची हमी देते. यावरून $f$ हे !!a{bijection} आहे.
$\Nat$ वर नेहमीचे आगमन वारंवार वापरून $f$ हे रचनारक्षक प्रतिचित्रण आहे हे
दाखवणे यापेक्षा अवघड नाही.

संचसैद्धान्तिक भाषेत फलन हा संबंधाचाच एक विशेष प्रकार असतो; उदाहरणार्थ,
$\lforall[x][\lexists![y][R(x,y)]]$ पूर्ण करणाऱ्या $R$ या द्विस्थानी
संबंधाशी एकस्थानी फलन~$f$ ची ओळख करता येते. त्यामुळे फलनांवरही संख्यापन करता
येते. मग पूर्ण चिन्हार्थमीमांसा वापरून अनंत !!{structure}s चा वर्ग असा
परिभाषित करता येतो: त्या !!{structure}s $\Struct M$ असतात ज्यांच्यासाठी
$\Struct M$ च्या !!{domain} पासून त्याच्याच उचित उपसंचाकडे जाणारे एक
एकास-एक फलन असते:
\[
\lexists[f][(\lforall[x][\lforall[y][(\eq[f(x)][f(y)] \lif \eq[x][y])]]
  \land \lexists[y][\lforall[x][\eq/[f(x)][y]]])].
\]
या वाक्याचा निषेध मग सांत !!{structure}s चा वर्ग परिभाषित करतो.

याखेरीज, रेषीय क्रमाच्या परिभाषेत पुढील अट जोडून सुक्रमांचा वर्ग परिभाषित
करता येतो:
\[
\lforall[P][(\lexists[x][\Atom{P}{x}] \lif \lexists[x][(\Atom{P}{x}
    \land \lforall[y][(y < x \lif \lnot \Atom{P}{y})])])].
\]
“संच” आणि “एकस्थानी संबंध” यांची ओळख केल्यास, प्रत्येक रिक्तेतर संचाला
लघुतम घटक असतो असे हे सांगते. आणखी एका उदाहरणात, शिखरांची परस्पर न जोडलेल्या
भागांत कोणतीही अ-तुच्छ विभागणी नसते असे सांगून आलेखांची संयोजकता व्यक्त करता
येते:
\[
\lnot \lexists[A][(\lexists[x][A(x)] \land \lexists[y][\lnot A(y)]
  \land \lforall[w][\lforall[z][((\Atom{A}{w} \land \lnot \Atom{A}{z})
      \lif \lnot \Atom{R}{w,z})]])].
\]
दुसऱ्या उदाहरणासाठी, ज्या सांत !!{structure}s च्या !!{domain} चा आकार सम
आहे त्यांचा वर्ग परिभाषित करण्याचा सराव तुम्ही करू शकता. आणखी लक्षवेधक
म्हणजे, परिमेय संख्या अंतर्भूत करणारे संपूर्ण क्रमित क्षेत्र म्हणून वास्तव
संख्यांचे समरूपतेपर्यंत एकमेव वर्णन देता येते.

थोडक्यात, द्वितीय-क्रम तर्कशास्त्र प्रथम-क्रम तर्कशास्त्रापेक्षा खूपच अधिक
अभिव्यक्तिसमर्थ आहे. ही चांगली बाजू; आता अडचण पाहू. पूर्ण द्वितीय-क्रम
चिन्हार्थमीमांसेसाठी संपूर्ण असणारी प्रभावी !!{derivation} प्रणाली नाही, हे
आपण आधीच नमूद केले आहे. प्रथम-क्रम तर्कशास्त्राचे अनेक गुणधर्म येथे उपलब्ध
नसतात; त्यांत संहतता आणि लोव्हेनहाइम--स्कोलेम प्रमेयांचा समावेश होतो.

दुसरीकडे, पूर्ण द्वितीय-क्रम चिन्हार्थमीमांसा सोडून दुर्बल चिन्हार्थमीमांसा
स्वीकारली, तर किमान द्वितीय-क्रम !!{derivation} प्रणाली या
चिन्हार्थमीमांसेसाठी संपूर्ण असते. दुसऱ्या शब्दांत, $\Proves$ चा अर्थ
“किमान प्रणालीत सिद्ध होते” आणि $\Entails$ चा अर्थ “दुर्बल
चिन्हार्थमीमांसेत तार्किकरीत्या निष्पन्न होते” असा घेतल्यास, $\Gamma \Entails
!A$ असेल तेव्हा $\Gamma \Proves !A$ असते, हे दाखवता येते.
!!{derivation} प्रणालीत ठरावीक गुणधर्माधारित संबंधरचना-स्वयंसिद्धके समाविष्ट
करायची असतील, तर ही स्वयंसिद्धके पूर्ण करणाऱ्या द्वितीय-क्रम रचनांपुरती
चिन्हार्थमीमांसा मर्यादित करावी लागते: उदाहरणार्थ, $\Delta$ हा
गुणधर्माधारित संबंधरचना-स्वयंसिद्धकांचा संच (शक्यतो त्यांपैकी सर्वांचा संच)
असेल, तर $\Gamma \cup \Delta \Entails !A$ असल्यास $\Gamma \cup \Delta
\Proves !A$ असते. विशेषतः, आपण विचारात घेत असलेली गुणधर्माधारित
संबंधरचना-स्वयंसिद्धके वापरून $!A$ सिद्ध करता येत नसेल, तर ती स्वयंसिद्धके
सत्य असूनही $\lnot !A$ सत्य असणारे एक प्रतिमान असते.

दुर्बल चिन्हार्थमीमांसेसाठी संपूर्णता प्रमेय का लागू होते हे पाहण्याचा सर्वांत
सोपा मार्ग म्हणजे द्वितीय-क्रम तर्कशास्त्राचा पुढीलप्रमाणे बहुवर्गीय
तर्कशास्त्र म्हणून विचार करणे. एका वर्गाचा अर्थ नेहमीचा “प्रथम-क्रम”
!!{domain} असा लावला जातो; मग प्रत्येक $k$ साठी “स्थानसंख्या $k$ असलेल्या
संबंधांचा” !!a{domain} असतो. भाषेत
“$\Atom{\Obj{true}_k}{R,x_1,\dots,x_k}$” ही अंगभूत संबंधचिन्हे घेतो. येथे
$R$ हे “$k$-स्थानी संबंध” या वर्गाचे चर आणि $x_1$, \dots,~$x_k$ या
प्रथम-क्रम वर्गातील वस्तू असताना, $R$ हा संबंध $x_1$, \dots,~$x_k$ यांना
लागू होतो असे हे चिन्ह सांगते.

ही ओळख स्वीकारल्यावर दुर्बल द्वितीय-क्रम चिन्हार्थमीमांसा मूलतः बहुवर्गीय
तर्कशास्त्राच्या नेहमीच्या चिन्हार्थमीमांसेसारखीच असते; आणि बहुवर्गीय
तर्कशास्त्राचे प्रथम-क्रम तर्कशास्त्रात अंतःस्थापन करता येते, हे आपण आधीच
पाहिले आहे. त्यामुळे दोन्ही दिशांतील रूपांतरणांचा विचार करता, द्वितीय-क्रम
तर्कशास्त्राची दुर्बल संकल्पना हे प्रत्यक्षात प्रथम-क्रम तर्कशास्त्राचेच
छद्मरूप आहे; त्यात !!{domain} मध्ये योग्य स्वयंसिद्धकांच्या अधीन असलेल्या
“वस्तू” आणि “संबंध” या दोन्हींचा समावेश असतो.

\end{document}
